% 四步仿射更新的关联扫描。上行：两两复合 g_{1:0}=g_1∘g_0，g_{3:2}=g_3∘g_2，根 g_{3:0}。
% 下行：根接收 id；左孩子继承，右孩子接收“左块总映射∘继承前缀”。
\begin{tikzpicture}[line width=.6pt,font=\fontsize{8}{10}\selectfont,>=Stealth,
  leaf/.style={draw=FigureBlue,fill=FigureBlueFill,rounded corners=1.8pt,minimum width=11mm,minimum height=6mm,inner sep=1pt},
  mid/.style={draw=FigureGreen,fill=FigureGreenFill,rounded corners=1.8pt,minimum width=17mm,minimum height=6mm,inner sep=1pt},
  pre/.style={draw=FigureYellow,fill=FigureYellowFill,rounded corners=1.8pt,minimum width=17mm,minimum height=6mm,inner sep=1pt},
  up/.style={->,draw=FigureStroke},
  down/.style={->,draw=FigureYellow,dashed},
  ttl/.style={FigureInk},
  note/.style={FigureMuted,font=\FigureNoteFont\fontsize{7.5}{9}\selectfont}]
% up-sweep
\node[ttl] at (25mm,25mm) {上行：合并相邻块};
\foreach \i/\x in {0/0,1/16,2/34,3/50}{\node[leaf] (l\i) at (\x mm,0) {$g_{\i}$};}
\node[mid] (m10) at (8mm,10mm) {$g_1\circ g_0$};
\node[mid] (m32) at (42mm,10mm) {$g_3\circ g_2$};
\node[mid,minimum width=22mm] (root) at (25mm,19mm) {$g_{3}\circ\cdots\circ g_0$};
\draw[up] (l0) -- (m10); \draw[up] (l1) -- (m10);
\draw[up] (l2) -- (m32); \draw[up] (l3) -- (m32);
\draw[up] (m10) -- (root); \draw[up] (m32) -- (root);
% down-sweep
\begin{scope}[xshift=72mm]
\node[ttl] at (25mm,25mm) {下行：传递前缀};
\node[pre,minimum width=11mm] (R) at (25mm,19mm) {$\mathrm{id}$};
\node[pre,minimum width=11mm] (L) at (8mm,10mm) {$\mathrm{id}$};
\node[pre] (Rr) at (42mm,10mm) {$g_1\circ g_0$};
\node[pre,minimum width=11mm] (p0) at (0,0) {$\mathrm{id}$};
\node[pre,minimum width=11mm] (p1) at (16mm,0) {$g_0$};
\node[pre] (p2) at (33mm,0) {$g_1\circ g_0$};
\node[pre,minimum width=21mm] (p3) at (53mm,0) {$g_2\circ g_1\circ g_0$};
\draw[down] (R) -- (L); \draw[down] (R) -- (Rr);
\draw[down] (L) -- (p0); \draw[down] (L) -- (p1);
\draw[down] (Rr) -- (p2); \draw[down] (Rr) -- (p3);
\node[note] at (25mm,-7mm) {叶子$i$再作用$g_i$，得到$\bm x_{i+1}=(g_i\circ p_i)(\bm x_0)$};
\end{scope}
\end{tikzpicture}
